\begin{table}[tb]
\centering\small\setlength{\tabcolsep}{3pt}
\caption{全24条件の資源・時間指標．Hのtoken欄は終端既知分の下限であり，不明２生成を０に補完していない．時間は実測合計．}
\begin{tabular}{lrr}
指標 & M & H \\ \hline
cacheを除く入力 & 276,458 & $\geq$296,062 \\
出力token & 66,022 & $\geq$91,970 \\
作業時間（秒） & 6,534.5 & 8,642.3 \\
準備・後片付け込（秒） & 6,570.7 & 8,679.0 \\
\hline
\end{tabular}
\end{table}
論理tokenの条件付き合計がHで小さくても，cacheを除く入力と出力は既知分だけでMの完全合計を上回る．増加はそれぞれ少なくとも7.09\%，39.30\%で，作業時間は実測32.26\%長かった．
入力処理と出力生成は異なる計算であり，token数をFLOPs，電力量，料金へ換算しない．時間にはモデル待ち，通知，計数，API等が含まれ，純粋なGPU処理時間ではない．ホスト負荷と実行経路も異なるため，時間差を通知方式の因果的な速度差とはしない．
両側成功部分ではcacheを除く入力が24,280，出力が9,537 token減り，作業時間も866.1秒短かった．一方，失敗部分は既知分だけでそれぞれ43,884，35,485 token増え，作業時間は2973.9秒長く，成功部分の減少を上回った．
P80/P81の別セッションの推論器滞在は計15,340.7秒，ロード20.12秒，評価40.42秒だった．これらは重なる測定値であり合算しない．共通時間をM/Hへ割り付けず，全件表に残す．
